% Generic AI workshop draft. Replace this layout with the venue's official
% style before submitting. Overleaf: main.tex, pdfLaTeX, BibTeX.
\documentclass[10pt,twocolumn]{article}
\usepackage[letterpaper,margin=0.75in,columnsep=0.25in]{geometry}
\usepackage[T1]{fontenc}
\usepackage[utf8]{inputenc}
\usepackage{lmodern}
\usepackage{microtype}
\usepackage{amsmath,amssymb,amsthm}
\usepackage{graphicx}
\usepackage{booktabs}
\usepackage[numbers,sort&compress]{natbib}
\usepackage[hidelinks]{hyperref}

% Change to \anonymousfalse for an attributed draft.
\newif\ifanonymous
\anonymoustrue
\newcommand{\papertitle}{Your Research Paper Title}
\newcommand{\methodname}{MethodName}
\newcommand{\R}{\mathbb{R}}
\newcommand{\E}{\mathbb{E}}
\newtheorem{theorem}{Theorem}
\newtheorem{proposition}{Proposition}
\newtheorem{definition}{Definition}

\title{\papertitle}
\ifanonymous
  \author{Anonymous Authors}
\else
  \author{First Author\thanks{Correspondence: author@example.com.}
    \quad Second Author\\Institution Name}
\fi
\date{}

\begin{document}
\maketitle

\begin{abstract}
% Aim for 150--200 words unless the selected workshop specifies otherwise.
Describe the problem and why it matters. State the limitation of existing
approaches, summarize the proposed idea, and report the main empirical or
theoretical finding once available. Close with the implication and scope
of the result. This paragraph is a writing prompt, not a research claim.
\end{abstract}

\section{Introduction}
\label{sec:introduction}
3D video games are a useful class of environments for assessing spacial intelligence in AI models.
\begin{itemize}
  \item
\end{itemize}
AI models are not very good at 3D video games. 

\paragraph{Contributions.}
% Replace these prompts with specific, defensible contributions.
\begin{itemize}
  \item A clearly defined research question or problem setting.
  \item A proposed method, analysis, or empirical observation.
  \item An evaluation that establishes the finding and its boundaries.
\end{itemize}

\section{Related Work}
\label{sec:related}
Group relevant literature by approach and explain how this work differs.
Use citations to support specific statements rather than listing papers.
Replace the example below with sources relevant to your research.

\paragraph{Citation example.}
\citet{vaswani2017attention} introduced the Transformer architecture.
A parenthetical citation looks like this: \citep{vaswani2017attention}.

\section{Method}
\label{sec:method}
\subsection{Problem Setup}
Define the inputs, outputs, assumptions, and evaluation objective.
For example, let $\mathcal{D}=\{(x_i,y_i)\}_{i=1}^{N}$ be a dataset and
let $f_\theta$ be a model with parameters $\theta$.

\subsection{Proposed Approach}
Describe \methodname{} precisely enough to reproduce it. Explain each
design choice and distinguish new components from established ones.
The following is an illustrative objective, to replace or remove:
\begin{equation}
  \label{eq:objective}
  \mathcal{L}(\theta)=\frac{1}{N}\sum_{i=1}^{N}
  \ell\bigl(f_\theta(x_i),y_i\bigr).
\end{equation}
Define the loss $\ell$ and any constraints when adapting
Equation~\eqref{eq:objective} to the actual method.

\begin{figure}[t]
  \centering
  % Replace the box with:
  % \includegraphics[width=\linewidth]{figures/overview.pdf}
  \fbox{\parbox[c][0.9in][c]{0.9\linewidth}{\centering
    Method overview figure\\Input $\rightarrow$ Model $\rightarrow$ Output}}
  \caption{Placeholder for a method overview. Explain each component
  and any symbols used in the final figure.}
  \label{fig:overview}
\end{figure}

\section{Experiments}
\label{sec:experiments}
\subsection{Experimental Setup}
% Include dataset provenance and licenses, split construction, leakage
% checks, baselines, metric definitions, seeds, and evaluation protocol.
Describe datasets, baselines, and metrics. State how hyperparameters
were selected and keep the test set separate from model selection.
Report model versions, hardware, and computational budget.

\subsection{Main Results}
Table~\ref{tab:results} provides a layout for results. Dashes are placeholders;
no experimental results are claimed. Specify whether uncertainty is a
standard deviation, standard error, or confidence interval, and report
the number of independent runs.

\begin{table}[t]
  \centering
  \caption{Example layout only. Replace metric names, units, and entries
  with measured results. Arrows indicate the preferred direction.}
  \label{tab:results}
  \begin{tabular}{lcc}
    \toprule
    Method & Score $\uparrow$ & Cost $\downarrow$ \\
    \midrule
    Baseline A & --- & --- \\
    Baseline B & --- & --- \\
    \methodname{} & --- & --- \\
    \bottomrule
  \end{tabular}
\end{table}

\subsection{Ablations and Analysis}
Isolate the effect of key components under comparable budgets.
Include representative failure cases and discuss variability or
negative findings that affect the interpretation of the results.

\section{Discussion and Limitations}
Describe what the evidence supports, where it does not generalize,
and which assumptions constrain the conclusions. Discuss relevant
privacy, bias, misuse, and societal impacts where applicable.

\section{Conclusion}
Restate the research question, supported finding, and practical or
scientific implication. Identify a concrete direction for future work.
Avoid introducing new results here.


\ifanonymous\else
\section*{Acknowledgments}
% Add contributors, funding, and disclosures for the attributed version.
\fi

\bibliographystyle{plainnat}
\bibliography{references}

\clearpage
\appendix
\section{Reproducibility Details}
\label{app:reproducibility}
% Record model checkpoints and versions, dataset versions, preprocessing,
% training settings, optimizer, schedule, seeds, hardware, runtime, and
% exact commands or configuration files needed to reproduce the results.
Provide implementation details that do not fit in the main paper.
Use anonymous artifact links when required by the selected workshop.

\section{Additional Results}
Include extended comparisons, sensitivity analyses, and failure cases.

\section{Proofs}
% Delete this section if the paper has no theoretical results.
Place supporting derivations and proofs here when applicable.

\end{document}
